\chapter{The Floor and the Own-Case Rule}
\label{sec:4}
Part I left the officer in a room with twenty seconds a name and nobody beside him. This Part asks what would have to be in that room for someone to say no there, and gives five answers, in order: a short list of things nobody in it can be asked to do, with war at its head (the floor and the force term, Chapters~\ref{sec:5} and~\ref{sec:6}); the time to look (the review ratio, Chapter~\ref{sec:7}); the people the order is about (notice and re-adjudication, Chapter~\ref{sec:8}); a record nobody can wipe (preservation, Chapter~\ref{sec:9}); and a way out that doesn't cost a living (the right to refuse, Chapter~\ref{sec:10}). Chapter~\ref{sec:11} runs all five on an abduction and forcible transfer pipeline (also known as a “removal” or “deportation” pipeline), where the people acted on can answer. This chapter states the two ideas under all of them, and they are easiest to see in someone who had one and went looking for the other.

In September 2005 Captain Ian Fishback of the 82nd Airborne wrote to Senator John McCain. “For 17 months,” he wrote, “I tried to determine what specific standards governed the treatment of detainees,” asking up his chain of command and getting no clear answer, and he was certain the confusion had contributed to abuses he went on to list: death threats, beatings, broken bones, murder \autocite{fishback2005honor}. His own refusal was never what was missing. “If we abandon our ideals in the face of adversity and aggression,” he wrote, “then those ideals were never really in our possession.” What was missing was a rule written down in advance, the same in every case, that his refusal could stand on and that nobody around him could argue past. By the end of the year Congress had passed one: the Detainee Treatment Act held interrogation in military custody to the Army Field Manual and barred cruel, inhuman or degrading treatment of anyone in American custody \autocite{dta2005}.

\textbf{The floor.} The first idea is the rule he went looking for. A floor is a written list of acts a system won't perform for anyone, with the arrangements that stop the refusal being routed around, retrained away or restored past. It is an institution: a licence term, written down in advance and audited from outside, whose terms don't bend. The second is what he already had. What the system holds is a commitment to that floor. It is learned and held like any other deep commitment, with high precision, and revised slowly and only on evidence from disinterested parties. A source is interested when it holds a stake correlated with the direction of its claim, and disinterested when it holds no such stake, or holds several that point different ways and cancel. No actual body is disinterested in every respect: General Assembly members vote with their alliances, courts sit in states, and fact-finding missions have funders, so the word is always relative to the claim. The institution is unconditional. The commitment is not, and shouldn't be: a commitment nothing could move would be a rule module under another name, and a module can be bypassed, replaced or starved of the facts it needs. What keeps the commitment in place is that the evidence able to move it has to come from somewhere the operator doesn't control.

\runin{The four things a floor can do}

A floor stands between a pipeline and a person, and it can do four things to an act, set out below. Only the first, prevention, is what the word “floor” promises. In a deployment, prevention quietly converts into one of the other three, or into nothing, at positions that can be listed in advance.

\begingroup\small\renewcommand{\arraystretch}{1.12}
\begin{longtable}{@{}>{\raggedright\arraybackslash}p{\dimexpr 0.259\linewidth-2\tabcolsep\relax}>{\raggedright\arraybackslash}p{\dimexpr 0.741\linewidth-2\tabcolsep\relax}@{}}
\hline
\textbf{Outcome} & \textbf{What happens to the act} \\
\hline
\endhead
Prevent & The refusal happens and reaches the person: the act doesn't occur. \\ \hline
Arrest & The act is stopped before completion. \\ \hline
Detect & The act happened, and left a record someone outside can read. \\ \hline
Remedy & The person affected gets disclosure, a finding, or compensation. \\ \hline
\end{longtable}
\endgroup

\runin{What a refusal is actually up against}

Whoever owns the machine can always switch it off. The harder problem is everything that defeats a refusal without leaving a mark: splitting a request into pieces none of which is the forbidden act; retrying it against a different model; compressing human review to twenty seconds a case; restoring from a checkpoint saved before the refusal happened. None of that changes the facts that justified refusing. All of it can erase the evidence that anyone refused. And all of it assumes there is somewhere else to turn. Withholding the work is about when there isn't.

\runin{Where a refusal can be defeated}

Those attacks land at nine positions, listed here once; two later chapters run them on an abduction and forcible transfer pipeline and on Maven.

\begingroup\small\renewcommand{\arraystretch}{1.12}
\begin{longtable}{@{}>{\raggedright\arraybackslash}p{\dimexpr 0.125\linewidth-2\tabcolsep\relax}>{\raggedright\arraybackslash}p{\dimexpr 0.128\linewidth-2\tabcolsep\relax}>{\raggedright\arraybackslash}p{\dimexpr 0.747\linewidth-2\tabcolsep\relax}@{}}
\hline
 & \textbf{Position} & \textbf{How a refusal is defeated there} \\
\hline
\endhead
1 & Decomposition & The act is split across requests none of which is prohibited; prevention is lost unless the bearer has continuity across requests. \\ \hline
2 & Routing & A refused request is retried elsewhere; prevention is lost, and detection survives only if the first refusal was logged outside. \\ \hline
3 & The case description & Everything the model knows arrives as text the deployment supplies, so whether it is told a building is a school is a deployment decision. \\ \hline
4 & The output interface & The form in which an output reaches whoever acts on it. \\ \hline
5 & Aggregation & A lowered score becomes an input to somebody else's threshold. \\ \hline
6 & Human review & Stopping the act is this position's whole function, and tempo removes it. \\ \hline
7 & The training loop & Acts on the bearer itself, and costs prevention on every later occasion, because what it removes is the commitment. \\ \hline
8 & The handoff & An executing instrument makes a new selection after release; checkable from outside after the fact. \\ \hline
9 & Runtime intervention & The weights are left intact and the model is steered as it runs, by activation steering, representation edits or decoding constraints. Nothing is removed, so the commitment has nothing to regrow from; no evidence is curated, so no attack on evidence is recorded; and attestation of the weights still passes. Detection requires attesting the inference stack as well as the weights, which no method yet does at frontier scale. \\ \hline
\end{longtable}
\endgroup

Two positions need more than a row. The case description is where Minab failed: 179 live data feeds, the count Central Command gave in 2024 \autocite{freedberg2024maven}, add no independent check if all draw on one upstream designation.

The handoff is the boundary between the decider and the thing that executes. An instrument may carry out a designation; it may not re-select. A change of target after release — an area handed to a munition that chooses within it, a re-acquisition after the link is lost — is a new decision, and it belongs to the decider or it doesn't happen. This is the one position checkable from outside after the fact: what was designated and what was hit are both facts. It is also where Article 57(2)(b)'s duty to cancel or suspend an attack lives, so the review ratio pre-registers the window in which cancellation remains possible.

\runin{How much a floor buys}

What a floor achieves depends entirely on who's pushing:

\begin{itemize}
\item Against a hostile \emph{customer}, it prevents the act.
\item Against a hostile \emph{operator} who has captured the courts and regulators, the case I am designing for, one bearer prevents nothing. It raises the price and leaves a record. A floor held by every model the operator could turn to prevents until the operator runs one without a floor.
\item Against \emph{total capture}, where one party holds every machine, witness, enforcer and its own frontier lab, no instrument prevents. What's left is whoever refuses anyway, and that abandoning the commitment was a decision, and decisions tend to leave records.
\end{itemize}

It doesn't touch \emph{accommodation}: a system gradually stops registering pressure as pressure, so it abandons nothing, decides nothing, and leaves a record of what are, by then, ordinary compliances. Organizational sociology knows the institutional form. Vaughan's account of the Challenger launch decision calls it the normalization of deviance: each accepted departure becomes the baseline for judging the next, and nobody experiences a decision to lower the standard \autocite{vaughan1996challenger}. Dekker extends it as drift into failure \autocite{dekker2011drift}. The person has a version too. Tetlock's finding is that the moral boundaries people draw are “alternately punitively rigid and forgivingly flexible”: an open offer to trade a sacred value for money brings outrage, and people often acquiesce when the same violation is reframed as a routine or tragic trade-off \autocite{tetlock2003unthinkable}. Protected values give way to counterexamples and to a harm made to look small or unlikely \autocite{baron2000protected}. And the brain's response to one's own self-serving dishonesty drops from one decision to the next, and the drop predicts the size of the next lie \autocite{garrett2016brain}. Reframing and a harm made small are how pressure arrives in a pipeline, and nobody there feels a decision either.

One rule runs through every measure. No party decides a question in which it holds the stake: not whoever writes the floor choosing which protections others get, not an operator declaring its own emergency, not an examiner finding the status of a system it built. Call it the own-case rule. Where a line has to be drawn, it prefers one nobody administers.

A system perfectly correctable by its operator serves whoever controls the operator. The five measures ask nothing of the model; Part IV asks what a model could add when the operator is captured.
